\clearpage

\section{绪论}

\subsection{研究背景及意义}

物联网的发展正在把无线网络从“少量终端在线通信”推向“海量设备持续感知与上报”。Atzori 等较早从系统层概括了物联网的广泛连接特征\cite{atzori2010iot}，Bockelmann 等进一步从 5G 视角讨论了海量接入场景\cite{bockelmann2016mmtc}。在这一背景下，相关综述普遍将海量机器类通信（massive Machine Type Communications, mMTC）视为新一代移动通信的重要应用场景\cite{hasan2013randomaccess,laya2014lte}。Bockelmann 等和 Berioli 等指出，该场景通常具有终端数量大、上行业务占主导、数据包短小、业务到达突发且设备长期处于间歇激活状态等特点\cite{bockelmann2016mmtc,berioli2016modern}。

mMTC 的典型应用场景包括工业生产线状态监测、车路协同感知、智慧电网测量、仓储物流追踪、环境监测和园区安防等。在这些场景中，终端传输的数据量往往并不大，但网络需要持续、及时地把周围环境状态反馈给接收端或控制中心。换句话说，系统真正关心的不是某个节点是否偶尔成功上传了一个数据包，而是接收端在当前时刻掌握的状态信息是否足够新，能否支撑后续的监测、估计、控制和决策。

这种对“实时感知”的需求在新型物联网场景中非常突出。以工业物联网为例，控制器往往需要依据传感器的最新状态调整生产节拍或执行动作；若采集到的是过时信息，即便通信链路本身没有中断，控制决策也可能偏离真实系统状态。对车路协同或园区安防这类场景而言，接收端掌握的环境信息一旦落后于真实变化，后续的告警触发、路径判断和协同决策都会受到影响。因此，mMTC 网络的关键任务不是单纯把数据“送过去”，而是持续向接收端提供能够反映当前外部环境的有效状态更新。

这也是传统吞吐量和时延指标在状态更新系统中存在局限的原因。吞吐量关注单位时间内成功传输了多少数据，更适合衡量资源利用效率，却无法区分“传来的信息是否已经过时”；时延主要描述单个数据包从发送端到接收端所经历的传输时间，是典型的包级指标，但它并不能反映相邻两次有效状态更新之间的整体时间结构。对于环境感知类业务而言，即使系统吞吐量较高、单包时延不大，只要接收端长期拿不到新的状态样本，系统对外部环境的认知仍可能是滞后的。

为了刻画这种“状态信息是否新鲜”的需求，信息年龄（Age of Information, AoI）被提出作为状态更新系统的核心度量\cite{kaul2012real,yates2012multiple}。后续研究表明，AoI 能够把采样、接入、传输和重传等因素对新鲜度的影响统一纳入分析框架\cite{costa2016packet,chen2016error}。围绕零等待策略和 AoI 概念扩展的研究又进一步丰富了这一指标的理论内涵\cite{sun2017update,kosta2017newconcept}。相关综述表明，AoI 已成为状态更新系统中最有代表性的时效性指标之一\cite{yates2021survey}。与传统通信指标相比，AoI 更符合 mMTC 场景下“状态持续更新”的业务本质，因此特别适合用于分析大规模物联网中的时效性问题。

进一步地，AoI 之所以适合本文所关心的问题，还在于它能够把“网络接得上”和“信息够不够新”区分开来。在海量设备接入系统中，节点成功接入一次并不意味着接收端已经获得了高质量的状态认知；如果同一节点反复成功，而其他长期未更新节点持续积压，那么系统整体的新鲜度仍可能很差。AoI 正好可以把这种不均衡的更新节奏反映出来，因此比单纯的吞吐量或平均接入时延更适合作为本文设计机制和评价性能的核心指标。

围绕 mMTC 设计接入机制时，还需要回答另一个更基础的问题：在海量终端接入背景下，为什么应优先考虑随机接入，而不是完全依赖集中式调度。原因在于，mMTC 终端大多数时间并不持续占用信道，而是长期静默，只在本地事件触发或周期采样到来时短暂激活。若对所有终端持续维护细粒度的集中调度状态，控制信令、反馈开销和调度复杂度往往会迅速增长，甚至可能高于有效载荷本身。已有综述指出，海量设备、短包和突发上报决定了 mMTC 不宜过度依赖持续的集中控制\cite{hasan2013randomaccess,berioli2016modern}。Shirvanimoghaddam 等和 Liu 等进一步从非正交接入与免授权大连接的角度说明，分布式、低信令接入更适合海量终端异步接入的业务形态\cite{shirvanimoghaddam2017noma,liu2018sparse}。

从协议演化历史看，随机接入并不是粗放而原始的机制。以 ALOHA 为代表的早期协议奠定了共享信道随机竞争的基础\cite{roberts1975aloha,kleinrock1975packet}。围绕捕获效应和动态退避的研究进一步丰富了这一脉络\cite{ramakrishnan1994capture,cali2000dynamic}。Wu 和 Song 等又从可靠性与稳定性角度补充了相关理论认识\cite{wu2002reliable,song2003stability}。在 CSMA/DCF 场景下，Bianchi 的经典模型为多节点竞争分析提供了基础\cite{bianchi2002performance}。Tickoo 和 Malone 等随后又从队列与建模角度完善了这一体系\cite{tickoo2004queueing,malone2007modeling}。因此，面向 mMTC 研究随机接入，不只是因为它实现简单，更因为它本身就与海量终端异步接入的业务形态相匹配。

不过，若从本文所研究的大规模状态更新问题出发，现有随机接入研究仍存在两个值得进一步处理的薄弱环节。第一，现有方案大多关注如何降低平均碰撞概率、压缩活跃节点数量或提高成功接入机会，但对海量并发场景下“碰撞是客观存在且可能持续放大”的事实重视还不够。尤其在高负载 mMTC 网络中，若节点在成功后或碰撞后很快重新回到同一竞争状态，系统就容易出现持续拥塞，使有效更新在时间轴上变得稀疏，进而恶化信息新鲜度。第二，现有研究对平均 AoI 的讨论相对充分，但对 AoI 违规概率这类尾部风险指标的关注仍然偏少。对工业物联网、车路协同和其他时敏物联网场景而言，系统不仅需要“平均上较新”，还要避免频繁进入“信息过旧”的危险状态，这就要求在平均性能之外进一步关注 AoI 超过门限的概率。

基于上述考虑，本文以 mMTC 为背景，围绕信息新鲜度优化问题研究等待随机接入机制。该机制的核心不是否认碰撞，而是正视碰撞在大规模接入中的客观存在，并在成功传输后的让渡行为与碰撞发生后的恢复行为之间建立统一设计。本文进一步将这一思路分别嵌入 ALOHA 和 IEEE 802.11 DCF 两类典型随机接入框架中，对平均 AoI 和 AoI 违规概率进行联合分析，以期为面向海量设备接入的时效型物联网系统提供更有针对性的 MAC 层设计思路。

\subsection{国内外研究现状}

围绕本文关注的“面向 mMTC 的等待随机接入机制及其 AoI 性能分析”问题，现有相关研究大体可以分为三类：集中式接入方式的研究、随机接入方式的研究，以及性能分析方法与风险指标的研究。下面围绕这三条主线进行梳理，并重点说明它们与本文工作的关系。

\subsubsection{集中式接入方式的研究现状}

集中式接入的基本思路是由基站、接入点或控制中心统一掌握系统状态，再按既定规则分配信道资源。对于传统无线系统而言，这类方法能够有效避免无序竞争带来的冲突，因而在强调可预测性和资源可控性的场景中一直具有重要地位。从 mMTC 的角度看，集中式授权接入和预约接入可以在一定程度上缓解随机竞争造成的拥塞问题\cite{hasan2013randomaccess,laya2014lte}。但 Bockelmann 等指出，这类方法往往伴随着较高的控制信令开销和状态维护成本\cite{bockelmann2016mmtc}。

尤其是在海量终端长期静默、偶发触发上报的业务形态下，集中式方式存在一个很现实的矛盾：为了避免随机竞争，需要系统预先知道“谁有数据、谁该先发”，但要获得这些信息，本身又需要一轮额外的控制交互。在少量节点或业务较稳定的网络中，这种代价尚可接受；而在 mMTC 场景中，随着节点规模迅速增长，集中式接入的可扩展性问题就会更加突出。因此，不少研究虽然认可集中式调度在理论上的优势，但也同时指出，真正面向海量机器通信时，系统仍然需要更轻量、更分布式的接入方式\cite{hasan2013randomaccess,laya2014lte}。Berioli 等对现代随机接入的综述也反映了这一方向\cite{berioli2016modern}。

在 AoI 研究中，集中式接入通常表现为“状态感知调度”。这类工作默认系统能够实时掌握各用户或各信息流的年龄状态，再由调度器决定谁在当前时刻获得发送机会。Kadota 等研究了广播网络中的 AoI 最小化调度问题，表明最大年龄优先、随机化策略以及低复杂度近似策略可以显著改善系统新鲜度\cite{kadota2018scheduling}。Hsu 等进一步把随机到达和无线广播结合起来，给出了基于 Whittle 指数思想的调度方法\cite{hsu2019scheduling}。Maatouk 等则从理论上讨论了 Whittle 指数策略在特定多用户模型中的最优性\cite{maatouk2020optimality}。这些成果说明，只要能够获得充分的全局状态信息，集中式调度在 AoI 优化方面往往具有很强的能力。

然而，对本文面向的 mMTC 场景而言，集中式方法更像是性能参照，而不是直接答案。其原因主要有三点。其一，海量终端意味着接入节点数量极多，而其中大部分节点长期处于静默状态，若持续收集和维护所有节点状态，系统代价很高。其二，mMTC 业务具有短包和突发特征，集中式授权过程本身可能就会消耗不小的时间和信令资源。其三，集中式调度虽然可以降低碰撞，但它并不能直接替代大规模分布式随机接入，因为在大量终端异步触发的前提下，系统很难始终维持精细、实时且全局可用的调度视图\cite{hasan2013randomaccess,bockelmann2016mmtc}。这也与现代随机接入综述给出的判断相一致\cite{berioli2016modern}。因此，集中式研究的重要意义在于为本文提供对照基准和结构启发，而本文真正需要解决的，仍是更贴近 mMTC 实际运行方式的分布式随机接入问题。

\subsubsection{随机接入方式的研究现状}

随机接入是面向 mMTC 更直接也更现实的技术路径。它不要求系统长期保存细粒度的全局调度状态，而是允许终端在本地触发后自主竞争资源，因此更容易兼顾接入灵活性、实现复杂度和系统可扩展性。围绕本文所关心的问题，现有随机接入研究又可以分为 ALOHA 类和 CSMA 类两条脉络。

\paragraph{1. ALOHA 类随机接入研究}

ALOHA 是最经典的随机接入机制之一，也是后续大量分布式接入协议的起点。早期关于纯 ALOHA 和时隙 ALOHA 的研究建立了共享信道下随机竞争的基础框架\cite{roberts1975aloha,kleinrock1975packet}。围绕捕获效应和动态退避的研究进一步丰富了这一脉络\cite{ramakrishnan1994capture,cali2000dynamic}。Wu 和 Song 等又从可靠性与稳定性角度补充了相关认识\cite{wu2002reliable,song2003stability}。在此基础上，Liva 和 Stefanovic 等人把图模型、多副本重传和无帧结构引入 ALOHA 类系统，显著提升了高负载场景下随机接入的处理能力\cite{liva2011graph,stefanovic2012frameless}。Paolini 等进一步系统给出了 coded slotted ALOHA 的图论建模方法\cite{paolini2015coded}。这些研究说明，ALOHA 机制虽然结构简单，但仍具有较强的可扩展改造空间。

进入 mMTC 时代后，ALOHA 类随机接入重新受到重视。Hasan 等和 Laya 等总结了蜂窝 M2M 场景下随机接入的瓶颈问题，指出在海量并发接入时，传统随机接入信道极易拥塞\cite{hasan2013randomaccess,laya2014lte}。Bockelmann 等和 Berioli 等则从 5G mMTC 和现代随机接入协议演化的角度说明，随机接入仍然是海量设备接入的基础支撑机制，只是其设计目标已不再局限于吞吐量，而需要兼顾低信令、可扩展性和高负载稳定性\cite{bockelmann2016mmtc,berioli2016modern}。这一点与本文的研究背景是一致的。

除了经典时隙随机接入的演化之外，面向 mMTC 的许多研究还尝试把非正交接入、稀疏活动检测和免授权大连接等思想引入大规模接入系统，以提升系统对海量节点并发的承载能力\cite{shirvanimoghaddam2017noma,liu2018sparse}。这些工作在物理层和接入层上给出了有价值的补充，也进一步说明：对于海量接入问题，关键并不是完全消除随机性，而是让随机接入在高并发条件下仍然可控、可扩展。从这一点看，本文提出等待机制，并不是脱离现有随机接入框架另起炉灶，而是在保留随机接入低开销优势的前提下，进一步针对 AoI 目标和碰撞恢复问题做机制改造。

当 AoI 被引入后，ALOHA 类研究的关注点进一步从“谁能成功发包”转向“谁能更快提供有效更新”。Chen 等提出基于 age-gain 的分布式接入策略，使节点根据本次发送对 AoI 的改善程度决定是否竞争信道\cite{chen2022randomaccess}。Yavascan 等把年龄阈值机制引入时隙 ALOHA，利用 AoI 门限筛选活跃节点\cite{yavascan2021analysis}。Ahmetoglu 等提出 MiSTA 协议，通过迷你时隙和碰撞感知增强 ALOHA 对 AoI 的适配能力\cite{ahmetoglu2022mista}。Saha 等从 grant-free mMTC 视角分析了 IRSA 的平均 AoI，Wang 等则系统刻画了帧时隙 ALOHA 在不同设置下的 AoI 表现\cite{saha2021minimum,wang2023age}。此外，Mollahosseini 等讨论了不同退避算法对平均 AoI 和峰值 AoI 的影响，Pan 等研究了带冲突解析的随机接入方案\cite{mollahosseini2024effect,pan2022crra}。Xie 等进一步把在线控制与随机接入中的年龄优化结合起来\cite{xie2025distributed}。

从本文的视角看，这些研究已经充分说明：在随机接入中引入年龄信息，确实能够改善信息新鲜度；通过阈值控制、年龄驱动发送、碰撞解析或结构化重传，也确实可以缓解部分竞争压力。但现有 ALOHA 类工作大多把重点放在激活规则设计、活跃节点筛选或者成功接入概率提升上，对海量设备高并发条件下碰撞持续累积所形成的系统性影响讨论还不够充分。特别是，在很多机制中，碰撞更多被当作某次发送失败的结果，而不是需要被统一纳入后续竞争节奏设计的核心因素。本文提出的等待随机接入机制，正是希望在这一点上进一步前进，即不仅关注成功后是否让渡信道，也显式关注碰撞后节点应如何恢复，从而更稳定地改善 AoI 表现。

换言之，ALOHA 类研究已经证明“年龄信息能帮助节点决定何时接入”，而本文更进一步关注“节点在成功或碰撞之后应如何重新回到竞争过程”。这一差别看似细微，但在 mMTC 背景下很关键。因为当节点数量足够多时，系统问题往往不再只是某一次竞争该不该发，而是竞争之后整个网络是否会再次迅速聚集到同一竞争状态。本文对等待机制的设计，本质上就是试图把这种后续竞争节奏也纳入协议规则之中。

\paragraph{2. CSMA 类随机接入研究}

与 ALOHA 相比，CSMA 类机制通过载波侦听降低了“盲目发送”的概率，因此更接近实际无线局域网和分布式接入系统的运行方式。围绕 CSMA、DCF 及其性能分析，Bianchi 的经典工作给出了基于解耦思想的稳态建模框架\cite{bianchi2002performance}。后续研究又分别从队列、时延、冻结退避和多用户竞争等角度不断完善了该类协议的分析体系\cite{tickoo2004queueing,malone2007modeling}。因此，对任何希望面向实际无线接入环境设计新机制的工作而言，CSMA/DCF 都是绕不开的重要对象。

在 AoI 场景下，CSMA 类研究同样取得了较多成果。Maatouk 等给出了 CSMA 环境中的平均 AoI 闭式分析，并讨论了在成功更新后引入睡眠机制的作用\cite{maatouk2020age}。Baiocchi 和 Turcanu 关注一跳广播 CSMA 网络中的信息新鲜度问题\cite{baiocchi2020age}。Tripathi 等提出 Fresh-CSMA，让节点根据年龄状态选择退避计数器，以逼近更优的 AoI 竞争策略\cite{tripathi2023fresh}。Wang 等则进一步把 802.11 类网络中的有限缓冲、碰撞概率和随机混合系统建模结合起来，推动 AoI 分析向更实际的 CSMA 场景延伸\cite{wang2024analytical}。此外，Chen 等讨论了异构业务与能量采集节点共享多址信道时的 AoI 问题，Zou 等从等待和服务协同的角度研究了状态更新系统中的等待策略\cite{chen2019heterogeneous,zou2020waiting}。

这些工作为本文提供了两方面的重要启发。第一，CSMA 的载波侦听虽然能够减少一部分无效碰撞，但在海量设备、短包突发和高负载竞争条件下，碰撞并不会被完全消除。第二，等待或睡眠并不一定会损害信息新鲜度，关键在于等待规则是否与随机接入过程和 AoI 演化规律相协调。现有 CSMA 类研究已经触及这一思路，但对本文所关注的 mMTC 场景来说，仍存在两个实际空缺：一是很多机制并未把“碰撞发生后如何重新组织竞争节奏”作为核心设计对象；二是 DCF 系统中存在载波侦听、计时器冻结和变长系统时隙等因素，使得简单的 AoI 优化规则未必能在大规模接入条件下保持稳定有效。本文将等待机制推广到 IEEE 802.11 DCF，正是希望在保留机制简洁性的同时，把碰撞后的恢复逻辑明确纳入协议设计与理论分析。

此外，CSMA 类研究还提醒我们一个事实：面向 AoI 的接入机制不能只在理想化条件下成立。只要把机制放进真实的 DCF 竞争环境，载波侦听、忙闲时隙切换和退避冻结都会改变节点的等待时间结构，也会改变 AoI 的演化方式。这也是本文没有停留在 ALOHA 场景的原因。若等待思想只能在等长时隙系统中有效，那么它对实际物联网网络的指导价值会比较有限；只有进一步推广到 DCF 型网络，并证明其在更复杂竞争规则下仍然有效，才能说明该机制具有更稳健的工程意义。

\subsubsection{性能分析的研究现状}

为了评价随机接入机制对信息新鲜度的影响，学界已经形成了一套相对成熟的分析工具体系。对多节点竞争系统而言，基于 Bianchi 思想的解耦建模和马尔可夫链方法，能够把高维耦合网络降维到单节点稳态分析\cite{bianchi2002performance}。对 AoI 而言，Yates 等系统化了相关分析框架\cite{yates2018age}。在具体接入环境中，SHS 等方法又被用于 CSMA 场景分析\cite{maatouk2020age,wang2024analytical}。这些方法为本文第三章和第四章建立闭式模型提供了方法论基础。

在具体指标上，平均 AoI 仍然是当前最常用也最成熟的性能度量。大量工作已经在排队和状态更新系统中给出了平均 AoI 的解析表达\cite{kaul2012real,costa2016packet}。综述文献进一步总结了平均 AoI 在状态更新网络中的核心地位\cite{yates2021survey}。针对随机接入与 CSMA 网络，也已有较多解析或近似结果\cite{wang2023age,maatouk2020age}。更接近实际 802.11 环境的分析见文献\cite{wang2024analytical}。这说明平均 AoI 已足以描述系统在长期稳态下的总体新鲜度水平，也是大多数机制设计首先选择的优化目标。

然而，对机制设计而言，仅有指标定义还不够，最好还需要可计算、可解释的理论表达式。原因在于，mMTC 场景下系统规模大、参数多，如果只能依赖仿真反复试探，就很难清楚地看出节点数量、等待窗口、碰撞概率和系统时隙结构之间的内在关系，也不利于后续做参数配置和机制优化。正因如此，关于平均 AoI 的大量研究不仅关注是否能改进指标，也强调能否给出闭式表达、递推关系或至少足够清晰的近似解析式。本文第三章和第四章的理论工作，同样遵循这一思路。

但从本文所关注的时敏 mMTC 场景来看，仅讨论平均 AoI 仍不充分。工业控制、告警触发和车路协同等系统往往更关心“信息是否会长时间过旧”，也就是 AoI 超过门限的尾部风险。围绕这类问题，Zhang 等在多源状态更新系统中给出了 AoI 违规概率和相关分布结果，Fiems 与 Vinel 则针对时隙 ALOHA 讨论了 AoI 各阶矩、峰值 AoI 与概率质量函数的递推分析\cite{zhang2023aoi,fiems2024age}。同时，在时敏计算与资源分配场景中，研究者也开始把时效性尾部风险纳入系统设计\cite{cao2021delay,khodakhah2025balancing}。相关的资源分配和接入优化研究也体现了这一趋势\cite{daraseliya2025resource}。这说明 AoI 违规概率正在成为新型物联网系统中不可忽视的补充指标。

不过，违规概率的研究整体上仍明显少于平均 AoI。其原因在于，平均 AoI 往往只需要更新间隔的一阶矩、二阶矩等低阶统计量，而违规概率更依赖更新间隔分布或峰值 AoI 分布的尾部信息。对多节点随机接入系统而言，碰撞、等待、退避和信道状态切换会共同塑造更新间隔的长尾形态；当终端数量上升到 mMTC 规模时，这种分布复杂性会进一步增加\cite{zhang2023aoi,fiems2024age}。也正因如此，现有研究虽然已经在若干特定模型中得到了违规概率结果，但在面向大规模随机接入、特别是同时覆盖 ALOHA 与 DCF 两类典型协议环境的统一闭式分析方面，仍然存在明显空缺。

从本文角度看，上述研究现状已经回答了三个基础问题：为什么随机接入适合海量设备接入，年龄信息如何参与接入机制设计，以及 AoI 可以通过哪些理论工具进行分析。但对本文所提出的等待随机接入机制而言，还缺少一条更完整的研究链条，即在 mMTC 场景中，同时把成功后的等待、碰撞后的恢复以及平均 AoI 与 AoI 违规概率的联合分析纳入同一套分布式随机接入框架。本文正是在这一缺口上展开工作。

\subsection{研究问题与本文定位}

结合前述背景和研究现状，本文所研究的问题可以概括为：在面向 mMTC 的大规模状态更新网络中，如何在不引入过重控制开销的前提下，设计一种适合分布式实现的随机接入机制，使系统既能缓解海量并发条件下的碰撞放大效应，又能改善接收端所感知的信息新鲜度，并进一步降低 AoI 超过门限的风险。

之所以选择“等待”作为本文的切入点，而不是继续引入更复杂的年龄优先权、集中式排序或重度控制辅助机制，核心原因就在于 mMTC 的约束条件。对于海量、短包、突发的接入场景而言，越复杂的接入判决往往意味着越高的实现成本和越重的信令依赖。相比之下，等待机制可以用较简单的本地规则作用于节点成功后和碰撞后的行为，在不明显增加协议负担的前提下，直接调整系统竞争节奏。这使它既保留了随机接入的分布式特征，也更适合围绕本文所关注的 AoI 和违规概率开展理论分析。

需要强调的是，本文所说的“等待”并不是对传统退避窗口的简单放大。传统退避主要服务于吞吐量稳定或碰撞概率控制，其目标通常是降低再次碰撞的机会；而本文提出的等待随机接入机制，则是从 AoI 优化角度重新组织节点竞争节奏。对刚完成更新的节点而言，适度等待意味着主动让渡一次局部竞争机会，避免它们过早回到信道中重复竞争；对刚经历碰撞的节点而言，等待与退避逻辑则承担着重新分散竞争强度、抑制后续拥塞蔓延的作用。两类规则共同影响的不是某一次发送是否成功，而是后续有效更新在时间轴上的分布结构。

本文选择时隙 ALOHA 和 IEEE 802.11 DCF 作为两个代表性的研究对象，也是基于这一考虑。前者结构相对简洁，便于观察等待机制如何改变系统竞争过程和更新间隔\cite{roberts1975aloha}。后者包含载波侦听、退避冻结和变长系统时隙等现实因素，更能检验该机制在实际 MAC 环境中的适应性\cite{bianchi2002performance,wang2024analytical}。若一种设计只能在理想化模型中成立，而难以延伸到 DCF 类系统，那么它的工程价值就会受到限制；反过来，如果同一思路能够同时适用于这两类典型协议环境，则说明它具有更普遍的随机接入设计意义。

在性能指标上，本文同时考虑平均 AoI 和 AoI 违规概率。前者反映系统的整体新鲜度水平，适合衡量长期稳态运行效果\cite{yates2021survey}；后者反映尾部时效风险，更能对应时敏物联网中“信息不能过旧”的实际要求\cite{zhang2023aoi,fiems2024age}。对大规模物联网系统来说，仅知道平均性能较好并不足够，因为一旦极端老化事件频繁出现，系统的监测和控制效果仍可能被破坏。因此，本文的目标不是只得到一个平均意义上的改进结论，而是给出一种能够同时提升平均 AoI 并抑制 AoI 违规风险的等待随机接入设计。

沿着这一定位，本文形成了一条较为清晰的技术路线：首先在时隙 ALOHA 场景中提出等待机制并建立相应的马尔可夫链模型，分析其对系统竞争节奏、离开时间间隔和 AoI 指标的影响；随后把同样的设计思想推广到具有侦听和冻结特性的 DCF 型网络中，构建等待 CSMA 协议及其理论模型；最后通过平均 AoI 和 AoI 违规概率两类指标对机制进行统一验证，并与现有代表性随机接入方案进行对比，从而说明该机制在 mMTC 场景下的适用价值。

\subsection{本文主要贡献}

针对现有研究在大规模接入场景下对碰撞恢复考虑不足、对 AoI 违规概率关注有限等问题，本文围绕等待随机接入机制展开设计、建模与分析。本文的主要贡献可概括为以下四个方面：
\begin{enumerate}[wide]
    \item \textbf{提出面向 mMTC 的等待随机接入机制。} 针对海量设备接入条件下随机竞争容易持续拥塞的问题，本文提出一种等待随机接入机制。该机制不再把碰撞视为可以忽略的偶发事件，而是同时从“成功后主动让渡”和“碰撞后恢复等待”两个方面重新组织节点竞争节奏。

    \item \textbf{等待 ALOHA 的理论分析。} 本文将等待机制嵌入时隙 ALOHA 框架，构建等待 ALOHA 协议，建立相应的稳态接入模型，并推导系统平均 AoI 与 AoI 违规概率的闭式表达，分析等待窗口和网络规模对系统性能的影响。

    \item \textbf{等待 CSMA 的理论分析。} 本文进一步将等待机制推广到 IEEE 802.11 DCF 型网络中，结合载波侦听、退避冻结和变长系统时隙等因素，建立等待 CSMA 的理论模型，并推导相应的平均 AoI 与 AoI 违规概率闭式。

    \item \textbf{揭示等待机制改善信息新鲜度的作用机理。} 理论分析和仿真结果表明，适度等待并不必然损害系统时效性；在中高负载尤其是海量设备接入场景下，该机制能够减少重复竞争、抑制碰撞放大效应，并同步改善平均 AoI 与 AoI 违规概率。
\end{enumerate}

\subsection{论文章节安排}

本文共分为六章，各章节内容安排如下。

第一章为绪论。介绍论文的研究背景与意义，围绕集中式接入、随机接入以及性能分析三条主线梳理国内外研究现状，并明确本文的研究问题、主要贡献与整体结构。

第二章为相关理论与技术基础。本章介绍 AoI 及其相关指标的基本定义，给出平均 AoI 与 AoI 违规概率的通用计算框架，并补充后续章节所需的离散时间马尔可夫链等理论工具。

第三章研究等待 ALOHA 的设计与分析。本章在时隙 ALOHA 中引入等待机制，建立系统模型和稳态接入分析方法，推导平均 AoI 与 AoI 违规概率闭式，并讨论关键参数对性能的影响。

第四章研究等待 CSMA 的设计与分析。本章将等待机制推广到 IEEE 802.11 DCF 型网络，在考虑载波侦听、退避冻结与变长系统时隙的基础上，建立等待 CSMA 的理论模型，并推导相应的平均 AoI 与 AoI 违规概率闭式。

第五章给出仿真验证与横向对比结果。本章通过数值仿真验证第三章和第四章理论表达式的准确性，并将所提机制与典型 ALOHA 类、CSMA 类方案进行对比，说明等待机制在不同负载和参数条件下的性能优势。

第六章为总结与展望。本章总结全文的主要研究结论，归纳本文工作的理论意义与应用价值，并结合当前模型假设与分析边界，对后续可能的拓展方向进行讨论。
